\begin{exercisechunk}
\item \textbf{From wavelet coefficient decay to an approximation rate.}\suppmark
\label[exercise]{ex:l15-wavelet-truncation}
\exerciseprofile{2}{2}{2}{\exproof\quad\excalculation\quad\exinterpretation}{%
Convert scale-wise wavelet coefficient bounds into an $m$-term
approximation rate.}
Let $(\psi_{j,k})$ be an orthonormal wavelet basis of $L^2([0,1])$.
Suppose that, at every level $j\geq0$, at most $C_0 2^j$ coefficients are
nonzero and
\[
 |\langle f,\psi_{j,k}\rangle_{L^2}|
 \leq C_1 2^{-j(r+1/2)},
 \qquad r>0.
\]
Let $f_J$ retain the coarse-scale terms and all wavelet coefficients from
levels $0,\ldots,J-1$.
\begin{enumerate}[(a),beginpenalty=10000]
\item Use orthonormality to prove
\[
 \|f-f_J\|_{L^2}^2
 \leq
 \frac{C_0C_1^2}{1-2^{-2r}}\,2^{-2rJ}.
\]
\item Suppose the number $m_J$ of retained coefficients satisfies
$m_J\asymp2^J$. Prove that
\[
 \|f-f_J\|_{L^2}=O(m_J^{-r}).
\]
\item State how \cref{prop:wavelet-coefficient-decay} supplies the hypothesis
of this exercise when the $r$th derivative is uniformly bounded on the
intervals appearing in that proposition.
Identify the step in parts (a)--(b) that converts local smoothness information
into a global $L^2$ approximation bound.
\end{enumerate}
\end{exercisechunk}
